\chapter{Quantum Spectral Curve and Hexagon Form Factors}
\label{ch:planar-n4-quantum-spectral-curve-hexagon}

\ConventionsLorentzian

The mirror TBA is exact but highly redundant.  The quantum spectral curve
compresses the same spectral information into a finite system of analytic
functions.  This chapter states the \(P\mu\)-system, explains its weak-coupling
reduction to Baxter equations, and records the separate hexagon framework for
three-point functions.

\section{Pmu System}

Let \(u\in\C\) be the spectral parameter.  The quantum spectral curve uses four
functions \(P_a(u)\), \(a=1,\ldots,4\), and an antisymmetric matrix
\(\mu_{ab}(u)=-\mu_{ba}(u)\).  The functions \(P_a\) have one short branch
cut, conventionally on \([-2g,2g]\), on their principal sheet.  Analytic
continuation through this cut is denoted by a tilde.  Introduce
\[
  P^a=\chi^{ab}P_b,
\]
where \(\chi^{ab}\) is the fixed antisymmetric tensor with nonzero entries
chosen so that \(P_aP^a=0\).

The \(P\mu\)-system is
\begin{align}
  \widetilde P_a(u)&=\mu_{ab}(u)P^b(u),
  \label{eq:planar-n4-pmu-one}\\
  \widetilde\mu_{ab}(u)-\mu_{ab}(u)
  &=
  P_a(u)\widetilde P_b(u)-P_b(u)\widetilde P_a(u),
  \label{eq:planar-n4-pmu-two}\\
  \mu_{ab}(u+\ii)&=\widetilde\mu_{ab}(u).
  \label{eq:planar-n4-pmu-period}
\end{align}
State quantum numbers enter through large-\(u\) asymptotics of \(P_a\) and
\(\mu_{ab}\).  In the \(SL(2)\) sector with twist \(J\), spin \(S\), and
dimension \(\Delta\), these asymptotics encode the charges
\[
  (J,S,\Delta)
\]
through powers and leading coefficients constrained by the algebraic
relations of the QSC.
With the \(SL(2)\) conventions used in the previous chapter, the large-\(u\)
powers are
\[
\begin{aligned}
  (P_1,P_2,P_3,P_4)
  &\sim
  (A_1u^{-J/2},A_2u^{-J/2-1},A_3u^{J/2},A_4u^{J/2-1}),\\
  (\mu_{12},\mu_{13},\mu_{14},\mu_{24},\mu_{34})
  &\sim
  (u^{\Delta-J},u^{\Delta+1},u^\Delta,u^{\Delta-1},u^{\Delta+J}).
\end{aligned}
\]
The leading coefficients are not arbitrary.  The QSC gluing and Pfaffian
constraint imply, for a state with \(S\) Lorentz-spin magnons,
\[
\begin{aligned}
  A_1A_4
  &=
  \ii\,
  \frac{((J+S-2)^2-\Delta^2)((J-S)^2-\Delta^2)}
       {16J(J-1)},\\
  A_2A_3
  &=
  \ii\,
  \frac{((J-S+2)^2-\Delta^2)((J+S)^2-\Delta^2)}
       {16J(J+1)}.
\end{aligned}
\]
These formulae are convention-sensitive because they depend on the shift
normalization \(f^{[\pm1]}(u)=f(u\pm\ii/2)\), the short cut
\([-2g,2g]\), and the choice of \(P\)-basis.

\begin{quotedtheorem}[Quantum spectral curve spectral problem]
\label{qthm:planar-n4-qsc-status}
Within the planar \(\mathcal N=4\) integrability framework, a physical
single-trace state is specified by QSC analyticity, gluing, asymptotic, and
regularity data.  Solving the resulting \(P\mu\)-system determines the exact
planar scaling dimension of the state as a function of \(\lambda\).
\end{quotedtheorem}

The theorem is quoted with its framework stated.  The QSC is a powerful exact
solution of the planar integrability system; it is not presently a
constructive definition of four-dimensional \(\mathcal N=4\) SYM.

\section{Weak-Coupling Baxter Limit}

At weak coupling the short cut collapses to the origin in the rescaled
spectral plane.  In the \(SL(2)\) sector, the QSC reduces to a finite-difference
Baxter equation for a polynomial \(Q(u)\) whose zeros are one-loop Bethe roots.
The equation has the form
\begin{equation}
  (u+\ii/2)^JQ(u+\ii)
  +(u-\ii/2)^JQ(u-\ii)
  =
  T(u)Q(u),
  \label{eq:planar-n4-qsc-baxter-limit}
\end{equation}
where \(T(u)\) is a degree-\(J\) transfer polynomial fixed by charges and
regularity.  If
\[
  Q(u)=\prod_{k=1}^S(u-u_k),
\]
then evaluating \eqref{eq:planar-n4-qsc-baxter-limit} at \(u=u_j\) gives
\[
  \left(\frac{u_j+\ii/2}{u_j-\ii/2}\right)^J
  =
  -\prod_{\substack{k=1\\k\neq j}}^S
  \frac{u_j-u_k-\ii}{u_j-u_k+\ii},
\]
which is the one-loop \(SL(2)\) Bethe equation in this convention.  Thus the
QSC contains the one-loop spin chain as its weak-coupling degeneration.

\begin{proposition}[Baxter equation from QSC degeneration]
\label{prop:qsc-baxter-degeneration}
Assume that in the weak-coupling limit the QSC functions in the \(SL(2)\)
sector reduce so that the ratio of two independent \(P\)-functions is a
degree-\(S\) polynomial \(Q(u)\), and that \(\mu\)-periodicity becomes a
finite-difference relation with shifts \(\pm\ii\).  Then regularity at the
zeros of \(Q\) gives the one-loop \(SL(2)\) Bethe equations.
\end{proposition}

\begin{proof}
Under the stated degeneration, the finite-difference relation is
\eqref{eq:planar-n4-qsc-baxter-limit}.  The right side is regular at every
zero \(u_j\) of \(Q\), so the left side must vanish there:
\[
  (u_j+\ii/2)^JQ(u_j+\ii)
  +(u_j-\ii/2)^JQ(u_j-\ii)=0.
\]
Dividing by \((u_j-\ii/2)^JQ(u_j-\ii)\) and using
\[
  \frac{Q(u_j+\ii)}{Q(u_j-\ii)}
  =
  -\prod_{\substack{k=1\\k\neq j}}^S
  \frac{u_j-u_k+\ii}{u_j-u_k-\ii}
\]
gives the displayed Bethe equation.
\end{proof}

\section{Cusp Anomalous Dimension and Konishi}

Two benchmark spectral quantities organize many tests of the framework.  The
first is the cusp anomalous dimension \(f(\lambda)\), extracted from large
spin twist operators:
\[
  \Delta-S
  =
  f(\lambda)\log S+O(S^0),
  \qquad S\to\infty.
\]
The QSC reproduces the same function obtained from the BES equation.  At weak
coupling it has a perturbative series in \(g^2\); at strong coupling it begins
with the semiclassical string result
\[
  f(\lambda)=\frac{\sqrt\lambda}{\pi}+O(\lambda^0)
\]
in the convention where \(\Delta-S=f(\lambda)\log S+O(S^0)\).

The second is the Konishi dimension.  The QSC treats it as an ordinary finite
state with \(J=2\), \(S=2\), and zero total momentum.  At weak coupling it
matches the asymptotic Bethe ansatz until wrapping corrections enter; at
strong coupling it gives the short-string expansion
\[
  \Delta_{\rm K}(\lambda)
  =
  2\lambda^{1/4}
  +O(\lambda^0).
\]
The numerical and analytic solution of the QSC is a calculation problem
inside the finite functional system \eqref{eq:planar-n4-pmu-one}--\eqref{eq:planar-n4-pmu-period}.

\section{Small Spin Expansion}

The QSC gives a controlled analytic expansion near \(S=0\), where the state
approaches the protected vacuum.  At leading order the \(\mu\)-matrix may be
chosen, up to QSC gauge transformations, so that
\[
  \mu_{12}=-1,\quad \mu_{13}=0,\quad
  \mu_{14}=\mu_{23}=1,\quad
  \mu_{24}=-\sinh(2\pi u),\quad \mu_{34}=0 .
\]
Writing
\[
  x(u)=\frac{u}{2g}+\sqrt{\frac{u^2}{4g^2}-1},
\]
with the short-cut branch, the leading solution has
\[
  P_1=\epsilon x^{-J/2},
  \qquad
  P_3=\epsilon(x^{-J/2}-x^{J/2}),
\]
and \(P_2,P_4\) determined by the Laurent expansion of \(\sinh(2\pi u)\) in
powers of \(x\).  Comparing the large-\(u\) coefficients with the algebraic
relations above gives
\[
  \Delta-J-S
  =
  \frac{4\pi g\,I_{J+1}(4\pi g)}{J\,I_J(4\pi g)}\,S+O(S^2),
\]
where \(I_J\) is the modified Bessel function.  Continuing the resulting
strong-coupling expansion to \(J=2,S=2\) gives the Konishi estimate
\[
  \Delta_{\rm K}
  =
  2\lambda^{1/4}+2\lambda^{-1/4}+O(\lambda^{-3/4}).
\]
The analytic continuation from small spin to integer spin is a QSC input, not
a theorem proved from the gauge-theory path integral.

\section{Hexagon Form Factors}

The spectral problem determines scaling dimensions.  Three-point functions
require additional data.  In the planar limit, a single-trace three-point
function can be represented by a pair of pants worldsheet.  The hexagon
program cuts this pair of pants into two hexagons and assigns a form factor to
each half.

For three single-trace operators with spin-chain lengths \(L_1,L_2,L_3\), the
bridge lengths are
\[
  \ell_{12}=\frac{L_1+L_2-L_3}{2},\qquad
  \ell_{23}=\frac{L_2+L_3-L_1}{2},\qquad
  \ell_{31}=\frac{L_3+L_1-L_2}{2}.
\]
Asymptotically, the structure constant is a sum over partitions of magnons
between the two hexagons:
\[
  C_{123}^{\rm asym}
  =
  \sum_{\alpha\cup\bar\alpha=\{u\}}
  (-1)^{|\bar\alpha|}
  \ee^{\ii p(\bar\alpha)\ell}
  H(\alpha)H(\bar\alpha)\,
  \mathcal N^{-1/2}.
\]
Here \(H\) is the hexagon form factor, \(p(\bar\alpha)\) is the total momentum
of the magnons moved across the bridge, and \(\mathcal N\) is the Gaudin norm
factor of the Bethe state.  Mirror particles on the bridges add finite-size
corrections, analogous in spirit to wrapping corrections in the spectral
problem.

\begin{openproblem}[Derivation of the hexagon expansion]
\label{op:planar-n4-hexagon-derivation}
Give a derivation of the hexagon expansion from a nonperturbatively defined
planar gauge-theory path integral, including the emergence of the hexagon
form factor, the measure for mirror corrections, and the treatment of
operator mixing at finite length.  Existing evidence is strong inside the
integrability framework.  A derivation at the level of a theorem in
four-dimensional QFT remains open.
\end{openproblem}

\section{Maldacena--Wilson Line, Cusp, and Bremsstrahlung}

Let \(C\) be an oriented smooth curve with parameter \(s\), tangent
\(\dot x^\mu(s)\), and scalar-coupling unit vector
\(\theta^I(s)\in S^5\).  The Maldacena--Wilson line in the fundamental
representation is
\[
  W_{\rm MW}(C,\theta)
  =
  \operatorname{Tr}\,P\exp\int_C
  \bigl(\ii A_\mu(x(s))\dot x^\mu(s)
  +\Phi_I(x(s))\theta^I(s)|\dot x(s)|\bigr)\,\dd s .
\]
It is a gauge-invariant line operator after taking the trace and imposing the
usual endpoint prescription for closed curves or external heavy sources.  A
straight line with constant \(\theta\) is half-BPS and has vanishing cusp
anomalous dimension.

For two rays meeting with Euclidean deflection angle \(\varphi\), define
\(\Gamma_{\rm cusp}(\varphi,\theta)\) by the logarithmic divergence of the
renormalized cusped line.  At small geometric angle and fixed scalar coupling,
\[
  \Gamma_{\rm cusp}(\varphi)
  =
  -B(\lambda)\varphi^2+O(\varphi^4).
\]
The coefficient \(B(\lambda)\) is the Bremsstrahlung function.  On the
straight-line defect, the displacement operator \(\mathbb D(t)\) has
\[
  \langle \mathbb D(t)\mathbb D(0)\rangle
  =
  \frac{C_D}{(t-\ii0)^4},
  \qquad C_D=12B(\lambda),
\]
and the radiated power of a slowly accelerating heavy probe is
\[
  \frac{\dd E_{\rm rad}}{\dd t}=2\pi B(\lambda)\,\dot v^2 .
\]
For planar \(\mathcal N=4\) SYM, supersymmetric localization of the circular
half-BPS Wilson loop gives the exact structural formula
\[
  B(\lambda)
  =
  \frac{\sqrt\lambda}{4\pi^2}
  \frac{I_2(\sqrt\lambda)}{I_1(\sqrt\lambda)}
  \qquad (N=\infty).
\]
This relation is a quoted exact input from the localization framework.  The
QSC and mirror-TBA descriptions of the cusped line reproduce the same
function with cusp-specific asymptotic data, providing a stringent
normalization check on \(g=\sqrt\lambda/(4\pi)\).

\section{Proof Boundary}

Planar integrability gives a highly constrained exact spectral structure for
four-dimensional interacting QFT.  Symmetry fixes matrix scattering, crossing
fixes the scalar factor, Bethe equations organize asymptotic spectra, mirror
TBA incorporates finite size, and the QSC compresses the exact spectral
problem into analytic functional equations.  The remaining mathematical task
is to derive this structure from a precise nonperturbative formulation of
planar \(\mathcal N=4\) SYM, including the topology in which local operators,
large-\(N\) limits, and analytic continuation in \(\lambda\) are controlled.
